\documentclass[11pt,a4paper]{article}

\usepackage[utf8]{inputenc}

\usepackage[T1]{fontenc}

\usepackage[polish]{babel}

\usepackage{lmodern}

\usepackage{geometry}

\geometry{margin=27mm}

\usepackage[hidelinks]{hyperref}

\title{Paradoks emocjonalnej adaptacji: ochrona przed dyskomfortem a tolerancja frustracji}

\author{Mateusz Pomianek}

\date{Szkic naukowy, wersja 1.0.0 | 28 listopada 2026}

\begin{document}

\maketitle

\begin{abstract}

Proponowana miara główna to prawdopodobieństwo poprawnego wznowienia zadania po wyłączeniu adaptacji.

\end{abstract}

\section{Wprowadzenie}

Model procesowy regulacji emocji rozróżnia modyfikację sytuacji i późniejsze reakcje na wymagania (Gross, 1998). Elastyczność regulacyjna uzasadnia ocenę zachowania po zmianie sytuacji (Bonanno \& Burton, 2013).

\section{Cel i pytania badawcze}

Hipotezy robocze wymagają zatwierdzenia przed rozpoczęciem właściwej weryfikacji.

\paragraph{H1.} Adaptacja redukująca trudność obniży deklarowaną frustrację podczas treningu względem adaptacji utrzymującej wyzwanie.

\paragraph{H2.} Po wyłączeniu adaptacji grupa utrzymująca wyzwanie częściej poprawnie wznowi zadanie w ciągu 30 s od standardowego zakłócenia niż grupa, w której uprzednio redukowano trudność.

\section{Materiał i metody}

Pierwszorzędowy wskaźnik: prawdopodobieństwo poprawnego wznowienia zadania po wyłączeniu adaptacji.

H2 stanowi kontrast główny w randomizowanym porównaniu. Zakłócenie i poziom kolejnej trudności muszą być identyczne we wszystkich grupach. Interfejs statyczny jest dodatkowym punktem odniesienia. H1 oceni odrębna skala bezpośrednio po treningu.

Przed badaniem z udziałem ludzi należy zatwierdzić wymagania etyczne, zasady zgody i ochronę danych.

\section{Plan analizy}

Zamrożenie estymandu, reguł oceny jakości i obsługi braków nastąpi przed dostępem do danych wynikowych. Raport obejmie oszacowanie kontrastu głównego z przedziałem niepewności i pełne mianowniki. Jeśli przewidziano hipotezy drugorzędowe, przed analizą należy ustalić sposób kontroli wielokrotności.

\section{Wyniki}

W niniejszym pliku nie umieszczono wyników empirycznych, statystyk ani nieprzeprowadzonych analiz. Sekcja wymaga uzupełnienia po uzyskaniu i weryfikacji materiału dowodowego.

\section{Dyskusja i ograniczenia}

Ograniczenia obejmują możliwe błędy pomiaru, dobór próby, alternatywne mechanizmy i ograniczoną możliwość uogólniania. Korelacja nie dowodzi przyczynowości; poprawność techniczna nie jest dowodem znaczenia klinicznego.

\section{Wnioski}

Hipotezy i sposób pomiaru stanowią podstawę oceny wykonalności oraz przygotowania właściwego protokołu.

\section*{Bibliografia}

{\setlength{\parindent}{0pt}
\setlength{\parskip}{5pt}
\hangindent=1.27cm
Bonanno, G. A., \& Burton, C. L. (2013). Regulatory flexibility: An individual differences perspective on coping and emotion regulation. \textit{Perspectives on Psychological Science, 8}(6), 591--612. \url{https://doi.org/10.1177/1745691613504116}\par
\hangindent=1.27cm
Gross, J. J. (1998). The emerging field of emotion regulation: An integrative review. \textit{Review of General Psychology, 2}(3), 271--299. \url{https://doi.org/10.1037/1089-2680.2.3.271}\par
}

\end{document}
